\uuid{5cOi}
\titre{ Calcul de l'accuracy pour un classifieur sigmoïde }
\niveau{L3}
\module{Analyse de données}
\chapitre{Réseaux de neurones}
\sousChapitre{Évaluation, métriques de classification}
\theme{Accuracy, seuil de décision, sigmoïde, classification binaire}
\auteur{Maxime Nguyen}
\datecreate{2026-03-19}
\organisation{AMSCC}
\difficulte{2}

\contenu{
\texte{
  On compare trois calculs d'accuracy pour un neurone de sortie sigmoïde produisant $\hat{y} \in [0, 1]$ :
  \begin{itemize}
    \item[A.] \texttt{preds = (y\_hat > 0.5).astype(int) ; acc = np.mean(preds == y)}
    \item[B.] \texttt{acc = np.mean(y\_hat == y)}
    \item[C.] \texttt{preds = (y\_hat > 0).astype(int) ; acc = np.mean(preds == y)}
  \end{itemize}
}
\begin{enumerate}

\item
  \question{Lequel est correct pour un neurone sigmoïde ?}
  \reponse{
    \textbf{A} est correct : on seuille $\hat{y}$ à $0{,}5$ pour obtenir une prédiction binaire $\in \{0, 1\}$,
    puis on compare à la vraie étiquette $y$.
  }

\item
  \question{Pourquoi B ne fonctionnera presque jamais ?}
  \reponse{
    La sortie sigmoïde $\hat{y}$ est un flottant dans $]0, 1[$ : la probabilité qu'il soit exactement
    égal à $y \in \{0, 1\}$ est quasi nulle. \texttt{np.mean(y\_hat == y)} retournera donc presque
    toujours $0$, ce qui n'a aucun sens comme mesure d'accuracy.
  }

\item
  \question{Dans quel cas C serait-il correct ?}
  \reponse{
    C seuille à $0$, ce qui serait pertinent si la sortie du réseau était une valeur brute non bornée
    (logit) plutôt qu'une probabilité sigmoïde. Par exemple, pour un neurone \texttt{tanh} dont la sortie
    est dans $]-1, 1[$, un seuil à $0$ permettrait de distinguer les deux classes.
    Pour une sigmoïde, le seuil naturel est $0{,}5$ (A), pas $0$ (C).
  }

\end{enumerate}
}
